\documentstyle[12pt]{article}

\newcommand{\SECT}[1]
		{\vspace{.3cm}
		$\!\!\!\!\!\!\!\!\!\!\!${\bf {#1}}
		\vspace{.2cm}}

\begin{document}


\begin{center}
	 {\bf Mathematics 477 \hspace{1em} Mathematical Theory of Probability
\par} 
\vspace{.5cm}

{\large \sc \bf Solutions to selected problems, I\par}
\vspace{.5cm}
{{\bf Yi-Zhi Huang} 
\par}
\end{center}

\paragraph{Chapter 1}

\begin{description}

\item[{\rm 7.}] 

\begin{description}

\item[(a)] Since there are 6 people, different ways correspond to 
different permutations of these people. So there are $6!$ ways.

\item[(b)] We can arrange them in the following way: First we decide
which group (the boys or the girls) should be the first three people in the
row. There are two ways.  Then we arrange the orders of the boys and the
girls. There are each 3! ways. So the answer is $2 \cdot 3! \cdot 3!$.

\item[(c)] We have four groups now: the group of $3$ boys and 
the three groups of three girls. We first arrange these four groups. There 
are $4!$ ways. Then we arrange the boys in the group. There are $3!$ ways.
The answer is $4!\cdot 3!$.

\item[(d)] We first decide whether a boy or a girl should be the first 
in the row. There are two ways. Then we arrange orders among the boys 
and among the girls. Each has $3!$ ways. So the answer is $2\cdot 3!\cdot 3!$.

\end{description}

\item[{\rm 28.}] 

\begin{description}

\item[(a)] Each teacher has four choices. So the total number of choices
is $4^{8}$ which is our answer.

\item[(b)] The number we want to calculate is the number of 
dividing $8$ objects into four groups of size $2, 2, 2, 2$. So the 
answer is ${8\choose 2, \; 2, \; 2, \; 2}$.

\end{description}

\end{description}

\paragraph{Chapter 2}

\begin{description}

\item[{\rm 13.}] For typographical reasons, we use $E, F, G$ to denote
$I, II, III$ in the problem. We think of the percentages as 
probabilities. Then the problem gives the following: 
\begin{eqnarray*}
&P(E)=0.1,\; \;\; P(F)=0.3,\;\;\; P(G)=0.05,&\\
&P(EF)=0.08,\;\;\;P(EG)=0.02,\;\;\; P(FG)=0.04,&\\
&P(EFG)=0.01.&
\end{eqnarray*}

\begin{description}

\item[(a)] We want to calculate first the percentage
$$P((EF^{c}G^{c})\cup (E^{c}FG^{c})\cup (E^{c}F^{c}G)).$$
Since $EF^{c}G^{c}$, $E^{c}FG^{c}$ and $E^{c}F^{c}G$ are 
mutually exclusive, we have 
\begin{eqnarray*}
\lefteqn{P((EF^{c}G^{c})\cup (E^{c}FG^{c})\cup (E^{c}F^{c}G))}\\
&&=P(EF^{c}G^{c})+P(E^{c}FG^{c})+P(E^{c}F^{c}G).
\end{eqnarray*}
Since $EF^{c}G^{c}=E(F\cup G)^{c}$ and 
$$(E(F\cup G)^{c})\cup (E(F\cup G))=E$$
and $E(F\cup G)^{c}$ and $E(F\cup G)$ are mutually exclusive,
we have $P(E)=P(E(F\cup G)^{c})+P(E(F\cup G))$. Thus
\begin{eqnarray*}
P(E(F\cup G)^{c})&=&P(E)-P(E(F\cup G))\\
&=&P(E)-P(EF\cup EG)\\
&=&P(E)-P(EF)-P(EG)+P(EFG)\\
&=&0.1-0.08-0.02+0.01\\
&=&0.01.
\end{eqnarray*}
Similarly, we obtain $P(E^{c}FG^{c})=0.19$ and $P(E^{c}F^{c}G)=0$.
Thus 
$$P((EF^{c}G^{c})\cup (E^{c}FG^{c})\cup (E^{c}F^{c}G))=0.20$$
and the answer is $20,000$. 

\item[(b)] We have 
\begin{eqnarray*}
\lefteqn{P((EF)\cup (EG)\cup (FG))}\\
&&=P(EF)+P(EG)+P(FG)\\
&&\quad -P(EFG)-P(EFG)-P(EFG)+P(EFG)\\
&&=P(EF)+P(EG)+P(FG)-2P(EFG)\\
&&=0.08+0.02+0.04-0.02\\
&&=0.12.
\end{eqnarray*}
So the answer is $12,000$.

\item[(c)] We have 
\begin{eqnarray*}
\lefteqn{P((EF)\cup (FG))}\\
&&=P(EF)+P(FG)-P(EFG)\\
&&=0.08+0.04-0.01\\
&&=0.11.
\end{eqnarray*}
So the answer is $11,000$.

\item[(d)] We have 
\begin{eqnarray*}
\lefteqn{P((E\cup F\cup G)^{c})}\\
&&=1-P(E\cup F\cup G)\\
&&=1-P(E)-P(F)-P(G)\\
&&\quad +P(EF)+P(EG)+P(FG)-P(EFG)\\
&&=1-0.1-0.3-0.05+0.08+0.02+0.04-0.01\\
&&=0.68.
\end{eqnarray*}
So the answer is $68,000$.

\item[(e)] We want to calculate $P((EFG^{c})\cup (E^{c}FG))$.
Note that $EFG^{c}$ and $E^{c}FG$ are mutually exclusive. 
Also note that since $(EFG^{c})\cup (EFG)=EF$ and $EFG^{c}$ 
and $EFG$ are mutually exclusive, $P(EF)=P(EFG^{c})+P(EFG)$,
or equivalently, $P(EFG^{c})=P(EF)-P(EFG)$. Similarly, 
$P(E^{c}FG)=P(FG)-P(EFG)$. Thus
\begin{eqnarray*}
\lefteqn{P((EFG^{c})\cup (E^{c}FG))}\\
&&=P(EFG^{c})+P(E^{c}FG)\\
&&=P(EF)-P(EFG)+P(FG)-P(EFG)\\
&&=0.08+0.01+0.04-0.01\\
&&=0.10.
\end{eqnarray*}
So the answer is $10,000$.


\end{description}

\item[{\rm 15.}] The number of elements of the sample space $S$ is 
${52\choose 5}$. 

\begin{description}

\item[(a)] We first select a suit. There are $4$ ways. Then 
we pick five cards from this suit. 
There are ${13\choose 5}$ ways. So the nunmber of elements in the 
event is $4\cdot {13\choose 5}$. The probability is 
$$\frac{4\cdot {13\choose 5}}{{52\choose 5}}.$$

\item[(b)] We first select four distinct kinds ($1$ or $J$
or $A$, etc.). There are 
${13\choose 4}$ ways. Then from these four kinds, we select one kind 
in which we shall pick two cards. There are $4$ ways. 
Finally we pick two cards from this kind and pick one card 
from each of the other three kinds. 
There are totally ${4\choose 2}\cdot 4\cdot 4\cdot 4$ ways. 
So the number of the elements of the event 
is ${13\choose 4} \cdot 4\cdot 
{4\choose 2}\cdot 4\cdot 4\cdot 4$. The probability is 
$$\frac{{13\choose 4} \cdot 4\cdot 
{4\choose 2}\cdot 4\cdot 4\cdot 4}{{52\choose 5}}.$$

\item[(c)] We first select three distinct kinds. There are 
${13\choose 3}$ ways. Then from these three kinds, we selct two 
kinds in which we shall pick two cards. There are ${3\choose 2}$ ways.
Finally we pick two cards from these two kinds and pick one card 
from the other kind. 
There are totally ${4\choose 2}\cdot {4\choose 2}\cdot 4$ ways. 
So the number of the elements of the event 
is ${13\choose 3} \cdot {3\choose 2}\cdot 
{4\choose 2}\cdot {4\choose 2}\cdot 4$. The probability is 
$$\frac{{13\choose 3} \cdot {3\choose 2}\cdot 
{4\choose 2}\cdot {4\choose 2}\cdot 4}{{52\choose 5}}.$$

\item[(d)] We first select three distinct kinds. There are 
${13\choose 3}$ ways. Then from these three kinds, we selct one
kind in which we shall pick three cards. There are $3$ ways.
Finally we pick three cards from this  kind and pick one card 
from each of the other kinds. 
There are totally ${4\choose 3}\cdot 4\cdot 4$ ways. 
So the number of the elements of the event 
is ${13\choose 3} \cdot 3\cdot 
{4\choose 3}\cdot 4\cdot 4$. The probability is 
$$\frac{{13\choose 3} \cdot 3\cdot 
{4\choose 3}\cdot 4\cdot 4}{{52\choose 5}}.$$

\item[(e)] We first select two distinct kinds. There are 
${13\choose 2}$ ways. Then from these two kinds, we selct one
kind in which we shall pick four cards. There are $2$ ways.
Finally we pick four cards from this  kind and pick one card 
from  the other kind. 
There are totally $4$ ways. 
So the number of the elements of the event 
is ${13\choose 2} \cdot 2\cdot  4$. The probability is 
$$\frac{{13\choose 2} \cdot 2\cdot  4}{{52\choose 5}}.$$


\end{description}

\item[{\rm 33.}] The assumption we make is that the chances of any 
elk (tagged or not tagged) being captured are equally likely. 
The number of the elements of the sample space is ${20\choose 4}$.
To have two elks tagged, we need to select two elks from the five 
tagged ones and select the other two from the fifteen untagged ones.
The number of such possible selections are ${5\choose 2}\cdot
{15\choose 2}$. So the probability is 
$$\frac{{5\choose 2}\cdot
{15\choose 2}}{{20\choose 4}}.$$


\item[{\rm TH 11.}] We prove the Bonferroni's inequality first:
\begin{eqnarray*}
P(EF)&=&1-P((EF)^{c})\\
&=&1-P(E^{c}\cup F^{c})\\
&=&1-P(E^{c})-P(F^{c})+P(E^{c}F^{c})\\
&=&1-(1-P(E))-(1-P(F))+P(E^{c}F^{c})\\
&=&P(E)+P(F)-1+P(E^{c}F^{c})\\
&\ge&P(E)+P(F)-1.
\end{eqnarray*}
In particular, when $P(E)=0.8$, $P(F)=0.9$,
$$P(EF)\ge P(E)+P(F)-1=0.8+0.9-1=0.7.$$

\end{description}

\paragraph{Chapter 3}

\begin{description}

\item[{\rm 32.}] Let $E$ be the event that the resigned person is a
woman and $F$ the event that the resigned person works in Store C. 
What we know is that 
\begin{eqnarray*}
P(E|F)&=&\frac{70}{100},\\
P(E|F^{c})&=&\frac{50\cdot 0.5+75\cdot 0.6}{50+75}\\
&=&\frac{70}{125},\\
P(F)&=&\frac{100}{50+75+100}\\
&=&\frac{100}{225},\\
P(F^{c})&=&1-P(F)\\
&=&1-\frac{100}{225}\\
&=&\frac{125}{225}.
\end{eqnarray*}
We want to calculate $P(F|E)$. By definition, 
$$P(F|E)=\frac{P(EF)}{P(E)}.$$
So we need to calculate $P(EF)$ and $P(E)$. 
First, 
$$P(EF)=P(F)P(E|F)=\frac{100}{225}\frac{70}{100}=\frac{70}{225}.$$
By Bayes' formula, we have 
\begin{eqnarray*}
P(E)&=&P(F)P(E|F)+P(F^{c})P(E|F^{c})\\
&=&\frac{100}{225}\frac{70}{100}+\frac{125}{225}\frac{70}{125}\\
&=&\frac{140}{225}.
\end{eqnarray*}
Thus 
$$P(F|E)=\frac{P(EF)}{P(E)}=\frac{1}{2}.$$

\item[{\rm 69.}] Let $E_{i}$ be the event that the $i$-th child is 
a boy. Then $E_{i}^{c}$ is the event that the $i$-th child is 
a girl. When $i\ne j$, $E_{i}$, $E_{i}^{c}$, $E_{j}$ and $E^{c}_{j}$ are 
all independent. Also we have $P(E_{i})=P(E_{i}^{c})=\frac{1}{2}.$

\begin{description}

\item[(a)] We have 
\begin{eqnarray*}
\lefteqn{P((E_{1}E_{2}E_{3}E_{4}E_{5})\cup (E_{1}^{c}E_{2}^{c}
E_{3}^{c}E_{4}^{c}E_{5}^{c}))}\\
&&=P(E_{1}E_{2}E_{3}E_{4}E_{5})+P(E_{1}^{c}E_{2}^{c}
E_{3}^{c}E_{4}^{c}E_{5}^{c})\\
&&=P(E_{1})P(E_{2})P(E_{3})P(E_{4})P(E_{5})\\
&&\quad +P(E_{1}^{c})P(E_{2}^{c})
P(E_{3}^{c})P(E_{4}^{c})P(E_{5}^{c})\\
&&=\frac{1}{2^{5}}+\frac{1}{2^{5}}\\
&&=\frac{1}{16}.
\end{eqnarray*}

\item[(b)] We have 
$$P(E_{1}E_{2}E_{3}E^{c}_{4}E^{c}_{5})=\frac{1}{2^{5}}=\frac{1}{32}.$$

\item[(c)] We first select where the boys should be (the oldest or the middle 
or the youngest, etc.) There are ${5\choose 3}$ ways. For each choice,
the probability is $\frac{1}{32}$. So the total probability 
is ${5\choose 3}\cdot \frac{1}{32}=\frac{5}{16}$. 

\item[(d)] $P(E_{1}^{c}E_{2}^{c})=\frac{1}{4}$.

\item[(e)] Let $F$ be the event of at least one fo the child is a girl.
Then $F^{c}$ is the event that none of the child is a girl. 
We know that $P(F^{c})=\frac{1}{32}$. So $P(F)=1-P(F^{c})=\frac{31}{32}$.

\end{description}

\item[{\rm 73.}] Note that for $A$ to be the match winner, $A$ must 
have won $k+2$ games if $B$ won $k$ games. Thus the total number of 
the games  played is $2k+2$ which is even. Similarly, if $B$ is the match 
winner, $B$ have won $k+2$ games if $A$ won $k$ games and the total number of 
the games  played is $2k+2$ which is also even. Note that 
we can group two consecutive games together. For $A$ be a match winner
when a total number of $2k+2$ games are played, 
the first $k$ groups must contain one game won by $A$ and one game won 
by $B$ and the last group must contain two games both won by $A$. 
Note that in the first $k$ groups, the orders whether $A$ or $B$ wins first
do not matter. Thus the number of elements of the 
event that $A$ wins at the $2k+2$ games is equal to 
$(2!)^{k}=2^{k}$. So the probability that $A$ wins at the $2k+2$ games
is $2^{k}p^{k+2}(1-p)^{k}$. Similarly the 
probability that $B$ wins at the $2k+2$ games
is $2^{k}p^{k}(1-p)^{k+2}$.


\begin{description}

\item[(a)] In this case, $k=1$. The probability is 
$2p^{3}(1-p)+2p(1-p)^{3}$.

\item[(b)] The probability that $A$ is the match winner is 
\begin{eqnarray*}
\sum_{k=0}^{\infty}2^{k}p^{k+2}(1-p)^{k}
&=&p^{2}\sum_{k=0}^{\infty}(2p(1-p))^{k}\\
&=&p^{2}\frac{1}{1-2p(1-p)}\\
&=&\frac{p^{2}}{1-2p+2p^{2}}.
\end{eqnarray*}

\end{description}




\end{description}



\end{document}